% Generated draft: requires mathematical and source review.
\documentclass[11pt]{article}
\usepackage[margin=25mm]{geometry}
\usepackage{fontspec}
\setmainfont{Latin Modern Roman}
\usepackage{amsmath,amssymb,mathtools,unicode-math}
\setmathfont{Latin Modern Math}
\usepackage{enumitem,needspace,xurl,hyperref}
\setlength{\parindent}{0pt}
\setlength{\parskip}{5pt}
\setlength{\emergencystretch}{3em}
\newcommand{\N}{\mathbb{N}}
\newcommand{\Z}{\mathbb{Z}}
\newcommand{\Q}{\mathbb{Q}}
\newcommand{\R}{\mathbb{R}}
\newcommand{\C}{\mathbb{C}}
\newcommand{\F}{\mathbb{F}}
\newcommand{\A}{\mathbb{A}}
\newcommand{\PP}{\mathbb{P}}
\newcommand{\E}{\mathbb{E}}
\newcommand{\eps}{\varepsilon}
\newcommand{\dd}{\,\mathrm{d}}
\begin{document}
\section*{UnsolvedMath ID: 30000253}
\textbf{Source problem code:} catalogue entry
\subsection{Definitions and mathematical statement}
The source requests a general determination of the number of spikes in steady states of the reaction--diffusion system labeled GM1. A steady state is a time-independent solution of that system, equivalently a solution of its stationary equations with its prescribed boundary conditions. Neither the GM1 equations nor those boundary conditions are reproduced in the supplied record.

The intended task is to characterize which spike counts and configurations occur among the relevant stationary solutions as the model data vary. If \(\mathcal S(\mathcal D)\) denotes the steady states for specified model data \(\mathcal D\), and \(s(U)\) is the source-defined spike count, the object to determine would be
\[
\{s(U):U\in\mathcal S(\mathcal D)\},
\]
possibly with asymptotic or geometric restrictions on the spike locations. This notation is only a description of the requested classification; the record supplies neither \(\mathcal D\) nor a definition of \(s\). Counting isolated local maxima, counting concentration centers in a singular limit, and counting components or higher-dimensional concentration sets are different notions. A single count independent of the chosen state or parameters is not asserted.

The supplied assessment checked 2026-08-21 reports later construction and counting of spike configurations in important geometries, including hyperplanes and line segments, but says that the unrestricted number-and-type question is not settled. It reproduces no model parameters, small-parameter regime, solution class, or exact count theorem. Thus no specific GM1 variant, diffusion scaling, or positivity assumption can be inserted merely from the system label.
\subsection{Short English statement}
Classify the possible spike counts of stationary GM1 solutions. The system equations, domain, boundary conditions, parameter regime, and definition of a spike are missing, so no unique numerical assertion can yet be formalized.
\subsection{Sources}
\begin{itemize}
\item Local UnsolvedMath registry, permanent ID 30000253, catalogue entry, \emph{Spike Counts in GM1 Reaction--Diffusion Steady States}.
\item \emph{Topological and Variational Methods for Differential Equations} (2006), Oberwolfach Reports, pp. 1601--1678, DOI \url{https://doi.org/10.4171/owr/2005/29}.
\item Zhang--Yang, supplied label ``number of spikes'': \url{https://doi.org/10.1080/00036811.2014.952491}.
\item Supplied literature triage checked 2026-08-21; no equations or precise hypotheses of the spike-count results are reproduced.
\end{itemize}
\end{document}
